
\documentclass[11pt,a4paper]{article}

% ------------------------------------------------------------
% Préambule commun minimal
% ------------------------------------------------------------
\usepackage[utf8]{inputenc}
\usepackage[T1]{fontenc}
\usepackage[french]{babel}
\usepackage{lmodern}
\usepackage{microtype}
\usepackage[a4paper,margin=1.65cm,headheight=15pt]{geometry}
\usepackage{amsmath,amssymb,mathtools}
\usepackage{array,tabularx,multicol}
\usepackage{enumitem}
\usepackage[most]{tcolorbox}
\usepackage{xcolor}
\usepackage{tikz}
\usepackage{pdflscape}
\usepackage{fancyhdr}
\usepackage{listings}
\usepackage{hyperref}

% ------------------------------------------------------------
% Couleurs et style
% ------------------------------------------------------------
\definecolor{fondpage}{RGB}{247,248,250}
\definecolor{encre}{RGB}{31,35,40}
\definecolor{bleu}{RGB}{40,91,135}
\definecolor{vert}{RGB}{55,125,92}
\definecolor{orange}{RGB}{178,105,42}
\definecolor{violet}{RGB}{101,77,145}
\definecolor{gris}{RGB}{86,91,99}
\definecolor{bleufonce}{RGB}{28,55,120}
\definecolor{bleuclair}{RGB}{238,243,255}
\definecolor{grisclair}{RGB}{245,245,245}

\hypersetup{
  colorlinks=true,
  linkcolor=bleufonce,
  urlcolor=bleufonce,
  pdftitle={Parcours de révision -- Probabilités -- Terminale spécialité},
  pdfauthor={}
}

\setlength{\parindent}{0pt}
\setlength{\parskip}{0.25em}
\renewcommand{\arraystretch}{1.12}
\setlength{\tabcolsep}{3pt}
\setlist[itemize]{leftmargin=1.25em,itemsep=0.15em,topsep=0.2em}
\setlist[enumerate,1]{label=\textbf{\arabic*.}, leftmargin=1.05cm, itemsep=0.3em, topsep=0.2em}
\setlist[enumerate,2]{label=\textbf{\alph*.}, leftmargin=0.85cm, itemsep=0.2em, topsep=0.15em}

\pagestyle{fancy}
\fancyhf{}
\lhead{}
\rhead{}
\cfoot{\thepage}

\newcommand{\Bin}{\mathcal B}
\newcommand{\origine}[1]{ {\scriptsize\textcolor{bleufonce}{(site #1)}}}


% Styles TikZ pour arbres de probabilités propres et lisibles
\tikzset{
  probaNoeud/.style={font=\small, inner sep=1.2pt},
  probaEtiq/.style={font=\footnotesize, fill=white, inner sep=1.35pt},
  probaBranche/.style={draw, line width=0.45pt, line cap=round},
  probaRacine/.style={circle, fill=black, inner sep=0.9pt}
}

% ------------------------------------------------------------
% Questions flashs : style repris du fichier source
% ------------------------------------------------------------
\newcounter{question}
\newenvironment{blocquestions}[1]{%
  \begin{tcolorbox}[
    enhanced,
    breakable,
    colback=bleuclair,
    colframe=bleufonce,
    boxrule=0.6pt,
    arc=2mm,
    left=2mm,right=2mm,top=2mm,bottom=1.5mm,
    title={\bfseries #1},
    coltitle=white,
    colbacktitle=bleufonce,
    before skip=3mm,
    after skip=3mm, fontupper=\large
  ]
  \large
  \begin{enumerate}[leftmargin=*,label=\textbf{\arabic*.},itemsep=0.45em,topsep=0.5em]
  \setcounter{enumi}{\value{question}}
}{%
  \setcounter{question}{\value{enumi}}
  \end{enumerate}
  \end{tcolorbox}
}

% ------------------------------------------------------------
% Carte mentale : macros reprises du fichier source
% ------------------------------------------------------------
\newcommand{\sep}{\vspace{0.55mm}}
\tcbset{
    enhanced,
    arc=2mm,
    outer arc=2mm,
    boxrule=0.42pt,
    left=1.15mm,
    right=1.15mm,
    top=0.75mm,
    bottom=0.75mm,
    boxsep=0.35mm,
    before skip=0.8mm,
    after skip=0.8mm,
    fonttitle=\bfseries\footnotesize,
    coltitle=encre,
    before upper=\raggedright\fontsize{8.05}{8.75}\selectfont,
}

\newtcolorbox{fichebox}[3][]{
    colback=white,
    colframe=#2!72!black,
    colbacktitle=#2!18,
    coltitle=#2!50!black,
    borderline west={0.9mm}{0pt}{#2!78!black},
    title={#3},
    drop fuzzy shadow=black!5,
    #1
}

\newcommand{\tagcol}[2]{%
\begin{tcolorbox}[
    enhanced,
    colback=#1!11,
    colframe=#1!45!black,
    boxrule=0.42pt,
    arc=2mm,
    left=1mm,right=1mm,top=0.45mm,bottom=0.45mm,
    before skip=0mm,
    after skip=0.85mm,
    fontupper=\bfseries\scriptsize,
    colupper=#1!55!black
]
\centering #2
\end{tcolorbox}}

\newtcolorbox{panel}[1]{
    colback=fondpage,
    colframe=fondpage,
    boxrule=0pt,
    arc=0mm,
    left=0mm,
    right=0mm,
    top=0mm,
    bottom=0mm,
    height=168mm,
    valign=top,
    before skip=0mm,
    after skip=0mm
}

\newcommand{\titrepage}{%
\begin{tcolorbox}[
    colback=white,
    colframe=bleu!70!black,
    boxrule=0.65pt,
    arc=3mm,
    left=2.5mm,right=2.5mm,top=1.15mm,bottom=1.15mm,
    before skip=0mm,
    after skip=1.4mm,
    drop fuzzy shadow=black!10
]
{\Large\bfseries Parcours de révision -- Probabilités}\hfill

\end{tcolorbox}}

% ------------------------------------------------------------
% Listings Python
% ------------------------------------------------------------
\lstset{
  language=Python,
  basicstyle=\ttfamily\small,
  numbers=left,
  numberstyle=\tiny,
  frame=single,
  showstringspaces=false,
  columns=fullflexible,
  keepspaces=true
}

\begin{document}

\begin{center}
{\LARGE\bfseries Parcours de révision -- Probabilités}\\[1mm]

\end{center}

\vspace{2mm}
\tableofcontents
\clearpage

% ------------------------------------------------------------
\phantomsection
\addcontentsline{toc}{section}{Partie 1 -- Récapitulatif de cours}
\newgeometry{margin=6.5mm}
\pagestyle{empty}
\begin{landscape}
\pagecolor{fondpage}\color{encre}
\thispagestyle{empty}
\fontsize{8.05}{8.75}\selectfont

\titrepage

\noindent
% ============================================================
% COLONNE 1
% ============================================================
\begin{minipage}[t]{0.242\linewidth}
\tagcol{bleu}{1. ÉVÉNEMENTS \& CONDITIONNEMENT}
\begin{panel}{bleu}

\begin{fichebox}{bleu}{Événements et notations}
\begin{itemize}
    \item $A$, $B$ : événements ; $\overline A$ : contraire de $A$.
    \item $A\cap B$ : $A$ \textbf{et} $B$ ; $A\cup B$ : $A$ \textbf{ou} $B$.
    \item $P_A(B)=P(B\mid A)$ : probabilité de $B$ \textbf{sachant} $A$.
\end{itemize}
\end{fichebox}

\begin{fichebox}{vert}{Probabilité conditionnelle}
Si $P(A)\neq 0$ :
\[
P_A(B)=\frac{P(A\cap B)}{P(A)}.
\]
Donc :
\[
P(A\cap B)=P(A)\times P_A(B).
\]
\textbf{Phrase :} $P_A(B)$ est la probabilité que $B$ soit réalisé sachant que $A$ est réalisé.
\end{fichebox}

\begin{fichebox}{vert}{Indépendance}
$A$ et $B$ sont indépendants si :
\[
P(A\cap B)=P(A)P(B).
\]
Si $P(A)\neq 0$ :
\[
P_A(B)=P(B).
\]
\textbf{Idée :} savoir que $A$ est réalisé ne change pas la probabilité de $B$.
\end{fichebox}

\end{panel}
\end{minipage}
\hfill
% ============================================================
% COLONNE 2
% ============================================================
\begin{minipage}[t]{0.242\linewidth}
\tagcol{orange}{2. ARBRES \& PROBABILITÉS TOTALES}
\begin{panel}{orange}

\begin{fichebox}{bleu}{Arbre pondéré}
\begin{center}
\resizebox{0.95\linewidth}{!}{%
\begin{tikzpicture}[x=1cm,y=1cm]
  \node[probaRacine] (O) at (0,0) {};
  \node[probaNoeud] (A) at (1.45,0.85) {$A$};
  \node[probaNoeud] (nA) at (1.45,-0.85) {$\overline A$};
  \node[probaNoeud] (AB) at (3.85,1.28) {$B$};
  \node[probaNoeud] (AnB) at (3.85,0.42) {$\overline B$};
  \node[probaNoeud] (nAB) at (3.85,-0.42) {$B$};
  \node[probaNoeud] (nAnB) at (3.85,-1.28) {$\overline B$};
  \draw[probaBranche] (O) -- (A.west) node[probaEtiq, midway, above, sloped] {$P(A)$};
  \draw[probaBranche] (O) -- (nA.west) node[probaEtiq, midway, below, sloped] {$P(\overline A)$};
  \draw[probaBranche] (A.east) -- (AB.west) node[probaEtiq, pos=.52, above, sloped] {$P_A(B)$};
  \draw[probaBranche] (A.east) -- (AnB.west) node[probaEtiq, pos=.52, below, sloped] {$P_A(\overline B)$};
  \draw[probaBranche] (nA.east) -- (nAB.west) node[probaEtiq, pos=.52, above, sloped] {$P_{\overline A}(B)$};
  \draw[probaBranche] (nA.east) -- (nAnB.west) node[probaEtiq, pos=.52, below, sloped] {$P_{\overline A}(\overline B)$};
\end{tikzpicture}%
}
\end{center}
\vspace{-1mm}
\begin{itemize}
    \item À un même nœud, la somme des branches vaut $1$.
    \item La probabilité d'un chemin est le \textbf{produit} des branches.
    \item La probabilité d'un événement est la \textbf{somme} des chemins qui le réalisent.
\end{itemize}
\[
P(A\cap B)=P(A)\times P_A(B).
\]
\end{fichebox}

\begin{fichebox}{violet}{Probabilités totales}
\textbf{Rédaction modèle :}

\medskip
\begin{center}
$A$ et $\overline A$ forment une partition de l'univers.

\smallskip
D'après la formule des\\ probabilités totales, on a :
\end{center}
\[
P(B)=P(A\cap B)+P(\overline A\cap B).
\]
Avec un arbre :
\[
\begin{aligned}
P(B)&=P(A)P_A(B)\\[-0.15em]
&\quad+P(\overline A)P_{\overline A}(B).
\end{aligned}
\]
\end{fichebox}

\end{panel}
\end{minipage}
\hfill
% ============================================================
% COLONNE 3
% ============================================================
\begin{minipage}[t]{0.242\linewidth}
\tagcol{vert}{3. VARIABLES ALÉATOIRES}
\begin{panel}{vert}

\begin{fichebox}{bleu}{Variable aléatoire}
Une variable aléatoire associe un nombre à chaque issue.

\medskip
\textbf{Phrase modèle :}
\[
\text{On note }X\text{ la variable aléatoire égale à }\ldots
\]
Toujours préciser ce que $X$ compte ou mesure.
\end{fichebox}

\begin{fichebox}{bleu}{Loi de probabilité}
Donner la loi de $X$, c'est donner toutes les valeurs possibles et leurs probabilités :
\[
\begin{array}{c|cccc}
x_i & x_1 & x_2 & \cdots & x_n\\
\hline
P(X=x_i)&p_1&p_2&\cdots&p_n
\end{array}
\]
avec :
\[
p_1+p_2+\cdots+p_n=1.
\]
\end{fichebox}

\begin{fichebox}{vert}{Espérance, variance, écart-type}
\[
E(X)=\sum x_i p_i
\]
\[
V(X)=\sum p_i\bigl(x_i-E(X)\bigr)^2
\]
\[
\sigma(X)=\sqrt{V(X)}.
\]
\textbf{Interprétation :} l'espérance est une \textbf{moyenne théorique} sur un très grand nombre de répétitions.
\end{fichebox}

\end{panel}
\end{minipage}
\hfill
% ============================================================
% COLONNE 4
% ============================================================
\begin{minipage}[t]{0.242\linewidth}
\tagcol{violet}{4. BERNOULLI \& LOI BINOMIALE}
\begin{panel}{violet}

\begin{fichebox}{bleu}{Bernoulli : une épreuve}
Une épreuve de Bernoulli a exactement deux issues :
\begin{itemize}
    \item succès $S$, de probabilité $p$ ;
    \item échec $\overline S$, de probabilité $1-p$.
\end{itemize}
\[
P(S)=p \qquad P(\overline S)=1-p.
\]
\end{fichebox}

\begin{fichebox}{violet}{Justifier une loi binomiale}
\textbf{Rédaction modèle :}

\smallskip
On considère l'épreuve de Bernoulli : \dotfill

\smallskip
Le succès est : « \ldots » ; sa probabilité est $p=\ldots$.

\smallskip
On répète cette épreuve $n$ fois de manière \textbf{identique et indépendante}.

\smallskip
$X$ compte le nombre de succès. Donc :
\[
X\sim\Bin(n;p).
\]
\end{fichebox}

\begin{fichebox}{vert}{Loi binomiale : formules}
Si $X\sim\Bin(n;p)$, alors, pour $0\leq k\leq n$ :
\[
P(X=k)=\binom nk p^k(1-p)^{n-k}.
\]
\[
E(X)=np\qquad V(X)=np(1-p)
\]
\[
\sigma(X)=\sqrt{np(1-p)}.
\]
\end{fichebox}

\begin{fichebox}{orange}{Traduire les événements}
\footnotesize
\begin{tabularx}{\linewidth}{>{\bfseries}X >{\centering\arraybackslash}X}
Expression & Probabilité \\
\hline
exactement $k$ & $P(X=k)$ \\
au plus $k$ & $P(X\leq k)$ \\
moins de $k$ & $P(X\leq k-1)$ \\
au moins $k$ & $P(X\geq k)$ \\
plus de $k$ & $P(X\geq k+1)$ \\
au moins un & $1-P(X=0)$
\end{tabularx}
\end{fichebox}

\end{panel}
\end{minipage}




\newpage
\thispagestyle{empty}

\begin{tcolorbox}[
    colback=white,
    colframe=bleu!70!black,
    boxrule=0.65pt,
    arc=3mm,
    left=2.5mm,right=2.5mm,top=1.15mm,bottom=1.15mm,
    before skip=0mm,
    after skip=1.4mm,
    drop fuzzy shadow=black!10
]
{\Large\bfseries Parcours de révision -- Probabilités}\hfill


\vspace{0.25mm}

\end{tcolorbox}

\noindent
% ============================================================
% PAGE 2 - COLONNE 1
% ============================================================
\begin{minipage}[t]{0.318\linewidth}
\tagcol{bleu}{5. CARDINAUX \& PRINCIPES DE COMPTAGE}
\begin{panel}{bleu}

\begin{fichebox}{bleu}{Cardinal et équiprobabilité}
Pour un ensemble fini $E$ :
\[
\operatorname{Card}(E)=\text{nombre d'éléments de }E.
\]
En équiprobabilité :
\[
P(A)=\frac{\operatorname{Card}(A)}{\operatorname{Card}(\Omega)}
=\frac{\text{cas favorables}}{\text{cas possibles}}.
\]
\textbf{Réflexe :} identifier l'univers $\Omega$, puis l'événement $A$.
\end{fichebox}

\begin{fichebox}{vert}{Réunion et complémentaire}
Pour deux ensembles finis :
\[
\begin{aligned}
\operatorname{Card}(E\cup F)
&=\operatorname{Card}(E)+\operatorname{Card}(F)\\[-0.2em]
&\quad-\operatorname{Card}(E\cap F).
\end{aligned}
\]
Si $E$ et $F$ sont disjoints :
\[
\operatorname{Card}(E\cup F)=\operatorname{Card}(E)+\operatorname{Card}(F).
\]
Si $A\subset E$ :
\[
\operatorname{Card}(\overline A)=\operatorname{Card}(E)-\operatorname{Card}(A).
\]
\end{fichebox}

\end{panel}
\end{minipage}
\hfill
% ============================================================
% PAGE 2 - COLONNE 2
% ============================================================
\begin{minipage}[t]{0.318\linewidth}
\tagcol{orange}{6. CHOIX SUCCESSIFS : L'ORDRE COMPTE}
\begin{panel}{orange}

\begin{fichebox}{vert}{Les trois formules}
Dans un ensemble à $n$ éléments :
\[
\begin{array}{c|c}
\text{Situation} & \text{Nombre}\\
\hline
k\text{ choix avec répétition} & n^k\\[0.2em]
k\text{ choix sans répétition} & A_n^k\\[0.2em]
\text{ordonner les }n\text{ éléments} & n!
\end{array}
\]
avec :
\[
A_n^k=n(n-1)\cdots(n-k+1)=\frac{n!}{(n-k)!}.
\]
\end{fichebox}

\begin{fichebox}{gris}{Exemples typiques}
\begin{itemize}
    \item Code à $k$ caractères : souvent $n^k$.
    \item Podium, ordre d'interrogation : arrangement.
    \item Ranger tous les objets : permutation $n!$.
\end{itemize}
\end{fichebox}

\end{panel}
\end{minipage}
\hfill
% ============================================================
% PAGE 2 - COLONNE 3
% ============================================================
\begin{minipage}[t]{0.318\linewidth}
\tagcol{vert}{7. CHOIX SANS ORDRE \& COEFFICIENTS BINOMIAUX}
\begin{panel}{vert}

\begin{fichebox}{vert}{Combinaisons et parties}
Choisir $k$ éléments parmi $n$, \textbf{sans tenir compte de l'ordre} :
\[
\binom nk=\frac{n!}{k!(n-k)!}.
\]
\emph{On choisit $k$ éléments parmi $n$, sans tenir compte de l'ordre. Il y a donc $\binom nk$ possibilités.}

\medskip
Nombre de parties d'un ensemble à $n$ éléments :
\[
2^n.
\]
Dans $\Bin(n;p)$, $\binom nk$ compte les placements des $k$ succès.
\end{fichebox}

\begin{fichebox}{vert}{Coefficients binomiaux utiles}
\[
\binom n0=1,\quad \binom nn=1,
\qquad \binom n1=n.
\]
\[
\binom nk=\binom n{n-k}.
\]
Relation de Pascal :
\[
\binom nk+\binom n{k+1}=\binom {n+1}{k+1}.
\]
Somme d'une ligne :
\[
\sum_{k=0}^{n}\binom nk=2^n.
\]
\end{fichebox}

\begin{fichebox}{orange}{Table de décision}
\begin{tabularx}{\linewidth}{>{\bfseries}X>{\centering\arraybackslash}X}
Question & Formule \\
\hline
ordre + répétitions & $n^k$ \\
ordre + éléments distincts & $A_n^k$ \\
ordonner tous les éléments & $n!$ \\
sans ordre & $\binom nk$ \\
toutes les parties & $2^n$
\end{tabularx}
\end{fichebox}

\end{panel}
\end{minipage}
\end{landscape}
\clearpage
\restoregeometry
\pagecolor{white}\color{black}
\pagestyle{fancy}

% ------------------------------------------------------------
\section*{Partie 2 -- Questions flashs}
\addcontentsline{toc}{section}{Partie 2 -- Questions flashs}
\setcounter{question}{0}
\begin{blocquestions}{A. Vocabulaire essentiel}
\item Qu'appelle-t-on l'univers en probabilités ? \origine{287}
\item Qu'appelle-t-on une expérience aléatoire ? \origine{310}
\item Qu'appelle-t-on une issue ? \origine{309}
\item Qu'appelle-t-on la loi de probabilité d'une variable aléatoire ? \origine{286}
\item Qu'est-ce qu'une variable aléatoire réelle ? \origine{75}
\item Que note-t-on souvent pour l'évènement contraire de \(A\) ? \origine{43}
\item Que signifie \(A\cap B\) ? \origine{41}
\item Que signifie \(A\cup B\) ? \origine{42}
\item Qu'appelle-t-on deux évènements incompatibles ? \origine{288}
\item Qu'appelle-t-on deux évènements indépendants ? \origine{289}
\item Qu'appelle-t-on une partition de l'univers ? \origine{312}
\item Qu'est-ce qu'une épreuve de Bernoulli ? \origine{57}
\end{blocquestions}

\begin{blocquestions}{B. Premiers calculs et arbres de probabilités}
\item Que vaut \(P(\overline{A})\) ? \origine{44}
\item Quelle est la probabilité de l'évènement impossible ? \origine{55}
\item Quelle est la probabilité de l'évènement certain ? \origine{56}
\item Si \(P(A)=0{,}35\), que vaut \(P(\overline{A})\) ? \origine{371}
\item Si \(A\) et \(B\) forment une partition de l'univers, que vaut \(P(A)+P(B)\) ? \origine{372}
\item Dans une épreuve de Bernoulli de probabilité de succès \(0{,}2\), quelle est la probabilité de l'échec ? \origine{373}
\item Quelle est la formule de la probabilité conditionnelle ? \origine{39}
\item Quelle formule permet de retrouver \(P(A\cap B)\) ? \origine{40}
\item Calculer \(P_B(A)\) sachant que \(P(A\cap B)=\dfrac15\) et \(P(B)=\dfrac3{10}\). \origine{660}
\item Comment calcule-t-on la probabilité d'un chemin dans un arbre ? \origine{46}
\item Comment calcule-t-on la probabilité d'un évènement obtenu par plusieurs chemins ? \origine{47}
\item Quelle est la formule des probabilités totales avec une partition \((B_i)\) ? \origine{49}
\end{blocquestions}

\begin{blocquestions}{C. Loi binomiale}
\item Qu'est-ce qu'un schéma de Bernoulli ? \origine{58}
\item Quand une variable aléatoire suit-elle une loi binomiale ? \origine{59}
\item Comment note-t-on une loi binomiale de paramètres \(n\) et \(p\) ? \origine{60}
\item Quelles valeurs peut prendre une variable \(X\sim\mathcal{B}(n,p)\) ? \origine{61}
\item Dans une loi binomiale, que compte la variable aléatoire ? \origine{221}
\item Quelle est la formule de \(P(X=k)\) pour \(X\sim\mathcal{B}(n,p)\) ? \origine{62}
\item Écrire \(P(X=3)\) si \(X\sim\mathcal{B}\!\left(8;\dfrac14\right)\). \origine{655}
\item Quelle est l'espérance d'une loi binomiale \(\mathcal{B}(n,p)\) ? \origine{63}
\item Que représente l'espérance dans un contexte concret ? \origine{64}
\item Si \(X\sim\mathcal{B}(10;0{,}2)\), que vaut \(E(X)\) ? \origine{222}
\item Quelle est la variance d'une loi binomiale \(\mathcal{B}(n,p)\) ? \origine{65}
\item Quel est l'écart-type d'une loi binomiale \(\mathcal{B}(n,p)\) ? \origine{66}
\item Quelle formule utilise-t-on souvent pour « au moins un succès » ? \origine{67}
\end{blocquestions}

\begin{blocquestions}{D. Variables aléatoires finies}
\item Donner la formule de l'espérance d'une variable aléatoire finie \(X\). \origine{489}
\item Donner la formule de la variance d'une variable aléatoire finie \(X\). \origine{490}
\item Comment définit-on l'écart-type d'une variable aléatoire \(X\) ? \origine{491}
\item Calculer l'espérance d'une variable aléatoire \(X\) telle que \(P(X=0)=\dfrac14\), \(P(X=1)=\dfrac12\), \(P(X=3)=\dfrac14\). \origine{665}
\item Quelle propriété vérifient les probabilités d'une loi finie ? \origine{540}
\item Quelle notation représente l'évènement « \(X\) prend la valeur \(a\) » ? \origine{542}
\item Comment calcule-t-on \(P(X\leqslant a)\) à partir de la loi de \(X\) ? \origine{543}
\end{blocquestions}

\begin{blocquestions}{E. Rédaction type bac}
\item Quels sont les points de rédaction importants pour justifier que \(X\) suit une loi binomiale ? \origine{396}
\item Quels sont les points de rédaction importants pour utiliser la formule des probabilités totales ? \origine{395}
\item On note \(X\) le nombre de succès obtenus en \(10\) essais indépendants. Un élève écrit : « \(X\) suit une loi binomiale car on répète \(10\) fois l'expérience. » Comment formuler proprement la réponse pour avoir tous les points ? \origine{335}
\item Pour trouver la probabilité d'un évènement \(B\) situé au bout des branches d'un arbre de probabilités, un élève écrit : « On applique la formule des probabilités totales avec \(A\) et \(\overline{A}\). » Comment formuler proprement la réponse pour avoir tous les points ? \origine{336}
\item Dans un exercice, un élève écrit : « La probabilité de \(A\cap B\), c'est \(P(A)\times P(B)\). » Comment formuler proprement la réponse pour avoir tous les points ? \origine{353}
\item À propos d'une loi binomiale, un élève écrit : « Au moins un succès, c'est \(1-P(X\geqslant 1)\). » Comment formuler proprement la réponse pour avoir tous les points ? \origine{414}
\end{blocquestions}

\clearpage

% ------------------------------------------------------------
\section*{Partie 3 -- Sujets de bac}
\addcontentsline{toc}{section}{Partie 3 -- Sujets de bac}
\section*{Sujet 1 -- Groupes sanguins et donneurs universels}
{\small\emph{Référence : Baccalauréat Métropole, 17 juin 2025, sujet 1, exercice 1.}}

On compte quatre groupes sanguins dans l'espèce humaine : A, B, AB et O. Chaque groupe sanguin peut présenter un facteur rhésus. Lorsqu'il est présent, on dit que le rhésus est positif, sinon on dit qu'il est négatif.

Au sein de la population française, on sait que :
\begin{itemize}
  \item 45 \% des individus appartiennent au groupe A, et parmi eux 85 \% sont de rhésus positif ;
  \item 10 \% des individus appartiennent au groupe B, et parmi eux 84 \% sont de rhésus positif ;
  \item 3 \% des individus appartiennent au groupe AB, et parmi eux 82 \% sont de rhésus positif.
\end{itemize}

On choisit au hasard une personne dans la population française.

On désigne par :
\begin{itemize}
  \item $A$ l'événement : « la personne choisie est de groupe sanguin A » ;
  \item $B$ l'événement : « la personne choisie est de groupe sanguin B » ;
  \item $AB$ l'événement : « la personne choisie est de groupe sanguin AB » ;
  \item $O$ l'événement : « la personne choisie est de groupe sanguin O » ;
  \item $R$ l'événement : « la personne choisie a un facteur rhésus positif ».
\end{itemize}

Pour un événement quelconque $E$, on note $\overline E$ l'événement contraire de $E$ et $p(E)$ la probabilité de $E$.

\begin{center}
\begin{tikzpicture}[x=1cm,y=1cm]
  \node[probaRacine] (racine) at (0,0) {};

  \node[probaNoeud] (A)  at (2.65, 1.80) {$A$};
  \node[probaNoeud] (B)  at (2.65, 0.60) {$B$};
  \node[probaNoeud] (AB) at (2.65,-0.60) {$AB$};
  \node[probaNoeud] (O)  at (2.65,-1.80) {$O$};

  \node[probaNoeud] (AR)  at (5.55, 2.12) {$R$};
  \node[probaNoeud] (AnR) at (5.55, 1.48) {$\overline R$};
  \node[probaNoeud] (BR)  at (5.55, 0.92) {$R$};
  \node[probaNoeud] (BnR) at (5.55, 0.28) {$\overline R$};
  \node[probaNoeud] (ABR)  at (5.55,-0.28) {$R$};
  \node[probaNoeud] (ABnR) at (5.55,-0.92) {$\overline R$};
  \node[probaNoeud] (OR)  at (5.55,-1.48) {$R$};
  \node[probaNoeud] (OnR) at (5.55,-2.12) {$\overline R$};

  \draw[probaBranche] (racine) -- (A.west)  node[probaEtiq, midway, above, sloped] {$\cdots$};
  \draw[probaBranche] (racine) -- (B.west)  node[probaEtiq, midway, above, sloped] {$\cdots$};
  \draw[probaBranche] (racine) -- (AB.west) node[probaEtiq, midway, below, sloped] {$\cdots$};
  \draw[probaBranche] (racine) -- (O.west)  node[probaEtiq, midway, below, sloped] {$\cdots$};

  \draw[probaBranche] (A.east) -- (AR.west)   node[probaEtiq, midway, above, sloped] {$\cdots$};
  \draw[probaBranche] (A.east) -- (AnR.west)  node[probaEtiq, midway, below, sloped] {$\cdots$};
  \draw[probaBranche] (B.east) -- (BR.west)   node[probaEtiq, midway, above, sloped] {$\cdots$};
  \draw[probaBranche] (B.east) -- (BnR.west)  node[probaEtiq, midway, below, sloped] {$\cdots$};
  \draw[probaBranche] (AB.east) -- (ABR.west) node[probaEtiq, midway, above, sloped] {$\cdots$};
  \draw[probaBranche] (AB.east) -- (ABnR.west)node[probaEtiq, midway, below, sloped] {$\cdots$};
  \draw[probaBranche] (O.east) -- (OR.west);
  \draw[probaBranche] (O.east) -- (OnR.west);
\end{tikzpicture}
\end{center}

\begin{enumerate}
  \item Recopier l'arbre ci-dessus en complétant les dix pointillés.
\item Montrer que $p(B\cap R)=0{,}084$. Interpréter ce résultat dans le contexte de l'exercice.
\item On précise que $p(R)=0{,}8397$. Montrer que $p_O(R)=0{,}83$.
\item On dit qu'un individu est « donneur universel » lorsque son sang peut être transfusé à toute personne sans risque d'incompatibilité. Le groupe O de rhésus négatif est le seul vérifiant cette caractéristique. Montrer que la probabilité qu'un individu choisi au hasard dans la population française soit donneur universel est de $0{,}0714$.
\item Lors d'une collecte de sang, on choisit un échantillon de 100 personnes dans la population d'une ville française. Cette population est suffisamment grande pour assimiler ce choix à un tirage avec remise. On note $X$ la variable aléatoire qui à chaque échantillon de 100 personnes associe le nombre de donneurs universels dans cet échantillon.
  \begin{enumerate}
    \item Justifier que $X$ suit une loi binomiale dont on précisera les paramètres.
    \item Déterminer à $10^{-3}$ près la probabilité qu'il y ait au plus 7 donneurs universels dans cet échantillon.
    \item Montrer que l'espérance $E(X)$ de la variable aléatoire $X$ est égale à $7{,}14$ et que sa variance $V(X)$ est égale à $6{,}63$ à $10^{-2}$ près.
  \end{enumerate}
\item Lors de la semaine nationale du don du sang, une collecte de sang est organisée dans $N$ villes françaises choisies au hasard numérotées $1,2,3,\ldots,N$, où $N$ est un entier naturel non nul.

  On considère la variable aléatoire $X_1$ qui à chaque échantillon de 100 personnes de la ville 1 associe le nombre de donneurs universels dans cet échantillon. On définit de la même manière les variables aléatoires $X_2$ pour la ville 2, ..., $X_N$ pour la ville $N$.

  On suppose que ces variables aléatoires sont indépendantes et qu'elles admettent la même espérance égale à $7{,}14$ et la même variance égale à $6{,}63$.

  On considère la variable aléatoire
  \[
  M_N=\dfrac{X_1+X_2+\cdots+X_N}{N}.
  \]
  \begin{enumerate}
    \item Que représente la variable aléatoire $M_N$ dans le contexte de l'exercice ?
    \item Calculer l'espérance $E(M_N)$.
    \item On désigne par $V(M_N)$ la variance de la variable aléatoire $M_N$. Montrer que
    \[
    V(M_N)=\dfrac{6{,}63}{N}.
    \]
    \item Déterminer la plus petite valeur de $N$ pour laquelle l'inégalité de Bienaymé-Tchebychev permet d'affirmer que
    \[
    P(7<M_N<7{,}28)>0{,}95.
    \]
  \end{enumerate}
\end{enumerate}

\newpage

\section*{Sujet 2 -- Chutes en roller et temps d'attente}
{\small\emph{Référence : Baccalauréat Métropole, 18 juin 2025, sujet 2, exercice 1.}}

Les deux parties peuvent être traitées indépendamment. Dans cet exercice on s'intéresse à des personnes venues séjourner dans un centre multisports au cours du week-end. Les résultats des probabilités demandées seront arrondis au millième si nécessaire.

\subsubsection*{Partie A}

Le centre propose aux personnes venues pour un week-end une formule d'initiation au roller composée de deux séances de cours.

On choisit au hasard une personne parmi celles ayant souscrit à cette formule.

On désigne par $A$ et $B$ les événements suivants :
\begin{itemize}
  \item $A$ : « la personne chute pendant la première séance » ;
  \item $B$ : « la personne chute pendant la deuxième séance ».
\end{itemize}

Pour un événement $E$ quelconque, on note $P(E)$ sa probabilité et $\overline E$ son événement contraire.

Des observations permettent d'admettre que $P(A)=0{,}6$.

De plus on constate que :
\begin{itemize}
  \item si la personne chute pendant la première séance, la probabilité qu'elle chute pendant la deuxième est de $0{,}3$ ;
  \item si la personne ne chute pas pendant la première séance, la probabilité qu'elle chute pendant la deuxième est de $0{,}4$.
\end{itemize}

\begin{enumerate}
  \item Représenter la situation par un arbre pondéré.
\item Calculer la probabilité $P(\overline A\cap \overline B)$ et interpréter le résultat.
\item Montrer que $P(B)=0{,}34$.
\item La personne ne chute pas pendant la deuxième séance de cours. Calculer la probabilité qu'elle n'ait pas chuté lors de la première séance.
\item On appelle $X$ la variable aléatoire qui, à chaque échantillon de 100 personnes ayant souscrit à la formule, associe le nombre d'entre elles n'ayant chuté ni lors de la première ni lors de la deuxième séance.

  On assimile le choix d'un échantillon de 100 personnes à un tirage avec remise.

  On admet que la probabilité qu'une personne ne chute ni lors de la première ni lors de la deuxième séance est de $0{,}24$.
  \begin{enumerate}
    \item Montrer que la variable aléatoire $X$ suit une loi binomiale dont on précisera les paramètres.
    \item Quelle est la probabilité d'avoir, dans un échantillon de 100 personnes ayant souscrit à la formule, au moins 20 personnes qui ne chutent ni lors de la première ni lors de la deuxième séance ?
    \item Calculer l'espérance $E(X)$ et interpréter le résultat dans le contexte de l'exercice.
  \end{enumerate}
\end{enumerate}

\subsubsection*{Partie B}

On choisit au hasard une personne venue un week-end au centre multisport. On note $T_1$ la variable aléatoire donnant son temps d'attente total en minute avant les accès aux activités sportives pendant la journée du samedi et $T_2$ la variable aléatoire donnant son temps d'attente total en minutes avant les accès aux activités sportives pendant la journée du dimanche.

On admet que :
\begin{itemize}
  \item $T_1$ suit une loi de probabilité d'espérance $E(T_1)=40$ et d'écart-type $\sigma(T_1)=10$ ;
  \item $T_2$ suit une loi de probabilité d'espérance $E(T_2)=60$ et d'écart-type $\sigma(T_2)=16$ ;
  \item les variables aléatoires $T_1$ et $T_2$ sont indépendantes.
\end{itemize}

On note $T$ la variable aléatoire donnant le temps total d'attente avant les accès aux activités sportives lors des deux jours, exprimé en minute. Ainsi on a $T=T_1+T_2$.

\begin{enumerate}
  \item Déterminer l'espérance $E(T)$ de la variable aléatoire $T$. Interpréter le résultat dans le contexte de l'exercice.
\item Montrer que la variance $V(T)$ de la variable aléatoire $T$ est égale à $356$.
\item À l'aide de l'inégalité de Bienaymé-Tchebychev, montrer que, pour une personne choisie au hasard parmi celles venues un week-end au centre multisports, la probabilité que son temps total d'attente $T$ soit strictement compris entre 60 et 140 minutes est supérieure à $0{,}77$.
\end{enumerate}

\newpage

\section*{Sujet 3 -- Jeu vidéo et suite probabiliste}
{\small\emph{Référence : Baccalauréat Asie, 11 juin 2024, sujet 2, exercice 2.}}

Léa passe une bonne partie de ses journées à jouer à un jeu vidéo et s'intéresse aux chances de victoire de ses prochaines parties.

Elle estime que si elle vient de gagner une partie, elle gagne la suivante dans 70 \% des cas.

Mais si elle vient de subir une défaite, d'après elle, la probabilité qu'elle gagne la suivante est de $0{,}2$.

De plus, elle pense avoir autant de chance de gagner la première partie que de la perdre.

On s'appuiera sur les affirmations de Léa pour répondre aux questions de cet exercice.

Pour tout entier naturel $n$ non nul, on définit les événements suivants :
\begin{itemize}
  \item $G_n$ : « Léa gagne la $n$-ième partie de la journée » ;
  \item $D_n$ : « Léa perd la $n$-ième partie de la journée ».
\end{itemize}

Pour tout entier naturel $n$ non nul, on note $g_n$ la probabilité de l'événement $G_n$. On a donc $g_1=0{,}5$.

\begin{enumerate}
  \item Quelle est la valeur de la probabilité conditionnelle $p_{G_1}(D_2)$ ?
\item Recopier et compléter l'arbre des probabilités ci-dessous qui modélise la situation pour les deux premières parties de la journée.
  \begin{center}
  \begin{tikzpicture}[x=1cm,y=1cm]
    \node[probaRacine] (r) at (0,0) {};
    \node[probaNoeud] (G1) at (2.15,0.92) {$G_1$};
    \node[probaNoeud] (D1) at (2.15,-0.92) {$D_1$};
    \node[probaNoeud] (G2a) at (4.85,1.28) {$G_2$};
    \node[probaNoeud] (D2a) at (4.85,0.56) {$D_2$};
    \node[probaNoeud] (G2b) at (4.85,-0.56) {$G_2$};
    \node[probaNoeud] (D2b) at (4.85,-1.28) {$D_2$};

    \draw[probaBranche] (r) -- (G1.west) node[probaEtiq, midway, above, sloped] {$\cdots$};
    \draw[probaBranche] (r) -- (D1.west) node[probaEtiq, midway, below, sloped] {$\cdots$};
    \draw[probaBranche] (G1.east) -- (G2a.west) node[probaEtiq, midway, above, sloped] {$\cdots$};
    \draw[probaBranche] (G1.east) -- (D2a.west) node[probaEtiq, midway, below, sloped] {$\cdots$};
    \draw[probaBranche] (D1.east) -- (G2b.west) node[probaEtiq, midway, above, sloped] {$\cdots$};
    \draw[probaBranche] (D1.east) -- (D2b.west) node[probaEtiq, midway, below, sloped] {$\cdots$};
  \end{tikzpicture}
  \end{center}
\item Calculer $g_2$.
\item Soit $n$ un entier naturel non nul.
  \begin{enumerate}
    \item Recopier et compléter l'arbre des probabilités ci-dessous qui modélise la situation pour les $n$-ième et $(n+1)$-ième parties de la journée.
    \begin{center}
    \begin{tikzpicture}[x=1cm,y=1cm]
      \node[probaRacine] (r) at (0,0) {};
      \node[probaNoeud] (Gn) at (2.15,0.92) {$G_n$};
      \node[probaNoeud] (Dn) at (2.15,-0.92) {$D_n$};
      \node[probaNoeud] (Gna) at (4.95,1.28) {$G_{n+1}$};
      \node[probaNoeud] (Dna) at (4.95,0.56) {$D_{n+1}$};
      \node[probaNoeud] (Gnb) at (4.95,-0.56) {$G_{n+1}$};
      \node[probaNoeud] (Dnb) at (4.95,-1.28) {$D_{n+1}$};

      \draw[probaBranche] (r) -- (Gn.west) node[probaEtiq, midway, above, sloped] {$g_n$};
      \draw[probaBranche] (r) -- (Dn.west) node[probaEtiq, midway, below, sloped] {$\cdots$};
      \draw[probaBranche] (Gn.east) -- (Gna.west) node[probaEtiq, midway, above, sloped] {$\cdots$};
      \draw[probaBranche] (Gn.east) -- (Dna.west) node[probaEtiq, midway, below, sloped] {$\cdots$};
      \draw[probaBranche] (Dn.east) -- (Gnb.west) node[probaEtiq, midway, above, sloped] {$\cdots$};
      \draw[probaBranche] (Dn.east) -- (Dnb.west) node[probaEtiq, midway, below, sloped] {$\cdots$};
    \end{tikzpicture}
    \end{center}
    \item Justifier que pour tout entier naturel $n$ non nul,
    \[
    g_{n+1}=0{,}5g_n+0{,}2.
    \]
  \end{enumerate}
\item Pour tout entier naturel $n$ non nul, on pose $v_n=g_n-0{,}4$.
  \begin{enumerate}
    \item Montrer que la suite $(v_n)$ est géométrique. On précisera son premier terme et sa raison.
    \item Montrer que, pour tout entier naturel $n$ non nul,
    \[
    g_n=0{,}1\times 0{,}5^{n-1}+0{,}4.
    \]
  \end{enumerate}
  \item Étudier les variations de la suite $(g_n)$.
\item Donner, en justifiant, la limite de la suite $(g_n)$. Interpréter le résultat dans le contexte de l'énoncé.
\item Déterminer, par le calcul, le plus petit entier $n$ tel que $g_n-0{,}4\leqslant 0{,}001$.
\item Recopier et compléter les lignes 4, 5 et 6 de la fonction suivante, écrite en langage Python, afin qu'elle renvoie le plus petit rang à partir duquel les termes de la suite $(g_n)$ sont tous inférieurs ou égaux à $0{,}4+e$, où $e$ est un nombre réel strictement positif.
\begin{lstlisting}
def seuil(e):
    g = 0.5
    n = 1
    while ...:
        g = 0.5 * ... + 0.2
        n = ...
    return n
\end{lstlisting}
\end{enumerate}

\end{document}
